\scanpage{262}

i.e. $(H,U_0\unsure{\in|\subset}H,\deleted{\unsure{\sigma\in H}}, T_H\subset H,\sigma\in H)$ satisfaisant les conditions

% slanted
\margin{$U_0,T_H$ ss-groupes de $H$}

\[\left\{\begin{aligned}
&\sigma^2=1\\
&U_0\cap U_1=\{1\}\qquad\text{où }U_1=\sigma(U_0)\\
&\sigma(t)=t^{-1}\qquad\forall t\in T_H\\
&T_H\text{ normalise }U_0\\
&U_0\times T_H\times U_1\unsure{-|\longrightarrow} H\qquad\text{injectif}\qquad(\text{où }U_1=\sigma(U_0)\\
&(u_0,t,u_1)\longmapsto u_0tu_1\hspace{12em}-
\end{aligned}\right.\tag*{(135)}\]

il reste à expliciter, en termes de ces données, le groupe $\widetilde Z$, les deux homomorphismes

\diagram{\begin{tikzcd}[ampersand replacement=\&]
{\deleted{\unsure{c,c'}}}\quad\widetilde Z\arrow[r,shift left,"c"] \& Z\arrow[l,shift left,"c'"]{\deleted{\unsure{=T_H}}}
\end{tikzcd}}

(où $Z$ est une copie conforme de $T_H$ via l'isom $u:Z\xrightarrow{\sim}T_H$)

et enfin l'opération de $\widetilde Z$ sur $H$, via

% slanted
\margin{donc c'est en fait l'opération de $T_0$ sur $H$}

\[\widetilde\lambda\longmapsto\operatorname{int}(u_0(\widetilde\lambda))|H\quad\deleted{\text{\unsure{écrire}}}\overset{\mathrm{déf}}{=}\Theta_0(\widetilde\lambda)\]

ou, au choix, via

% slanted
\margin{i.e. l'opération de $T_1$ sur $H$.}

\[\widetilde\lambda\longmapsto\operatorname{int}(u_1(\widetilde\lambda))|H\overset{\mathrm{déf}}{=}\Theta_1(\widetilde\lambda)\]

ces deux opérations étant \inserted{en effet} reliées par la formule -- déduite de (119)

\[u_1(\widetilde\lambda)=u_0(\widetilde\lambda)u(c(\widetilde\lambda))\]

qui donne

\[\Theta_1(\widetilde\lambda)=\Theta_0(\widetilde\lambda)\operatorname{int}(u(c(\widetilde\lambda))),-\tag*{(136)}\]

\deleted{Cela} de sorte que $\Theta_0(\widetilde\lambda),\Theta_1(\widetilde\lambda)$ se déterminent mutuellement. Il revient aussi au $=$ de \deleted{\unsure{de de|de}} demander, \inserted{plus symétriquement,} la construction du groupe $T_G$, comme groupe opérant sur $H$, avec dans $T_G$ les \deleted{deux \unsure{sous}-}\inserted{trois sous-}groupes $T_H\hookrightarrow T_G$ (plongement à déterminer), $T_0$ et $T_1$, \deleted{\unsure{ce qui est|que}} \inserted{$\widetilde Z$ étant} alors défini \deleted{\unsure{comme}} \inserted{de façon} \unsure{symétrique|symétriques} comme étant $T_G/T_H$, et l'isomorphisme $Z\xrightarrow{c'}\widetilde Z$ étant décrit comme le composé de

% archive sheet number: 261
